\documentclass[10pt]{beamer}
%\documentclass[handout,10pt]{beamer}
\usepackage{my_style}
%%%%%%%%%%%%%%%%%%%%%%%%%%%%%%%%%%%%%%%%%%%%%%%%%%%%%%%%%%%%%%%%%%%%
%%%%%%%%%%%%%%%%%%%%%%%%%%%%%%%%%%%%%%%%%%%%%%%%%%%%%%%%%%%%%%%%%%%%

\title[Math/CS 375]{Lecture 06 - A}
\subtitle{\S 2.1, 2.2 Gaussian Elimination}

\begin{document}

\begin{frame}
  \titlepage
\end{frame}

\begin{frame}{Goals}
We will answer these questions:\\[3mm]
\begin{itemize}
    \item How do you solve a \emph{general} linear system?
    \item How expensive is it to perform Gaussian elimination?
    \item What are some problems with the "naive" algorithm?
\end{itemize}
\end{frame}

\begin{frame}{Special Systems}
    \begin{itemize}
    \item \textbf{Diagonal} - Costs $\mO(n)$ FLOPs
    $
    \begin{bmatrix}
    \bullet & 0 & 0 & 0 \\
    0 & \bullet & 0 & 0 \\
    0 & 0 & \bullet & 0 \\
    0 & 0 & 0 & \bullet
    \end{bmatrix}
    $ 

    \item \textbf{Banded} - Costs $\mO(m^2n)$ FLOPs
    $
    \begin{bmatrix}
    \bullet & \bullet & 0 & 0 \\
    \bullet & \bullet & \bullet & 0 \\
    0 & \bullet & \bullet & \bullet \\
    0 & 0 & \bullet & \bullet
    \end{bmatrix}
    $
    
    \item \textbf{Triangular} - Costs $\mO(n^2)$ FLOPs
    $
    \begin{bmatrix}
    \bullet & \bullet & \bullet & \bullet \\
    0 & \bullet & \bullet & \bullet \\
    0 & 0 & \bullet & \bullet \\
    0 & 0 & 0 & \bullet
    \end{bmatrix}
    $

\end{itemize}
\question{
    \begin{itemize}
        \item What if $A$ is a general matrix? 
        \item What is the cost of solving $Ax=b$?
    \end{itemize}
}
\end{frame}


%%%%%%%%%%%%%%%%%%%%%%%%%%%%%%%%%%%%%%%%%%%%%%%%%%%%%%%%%%%%%%%%%%%%%%%%%
\sectionframe{Gaussian Elimination}
%%%%%%%%%%%%%%%%%%%%%%%%%%%%%%%%%%%%%%%%%%%%%%%%%%%%%%%%%%%%%%%%%%%%%%%%


% \begin{frame}{Sparse vs. Dense}
  
%     \begin{block}{Sparse Matrix}
%         \begin{itemize}
%             \item A matrix is considered \textbf{sparse} if the majority of its elements are zero. 
%             \item Example: triangular matrix, diagonal matrix, banded matrix 
%             \item $A$ is sparse with special pattern $\implies$ $Ax=b$ can be solved quickly 
%         \end{itemize}        
%     \end{block}
%         \begin{itemize}
%             \item Diagonal systems cost $\mO(n)$ FLOPs
%             \item Triangular systems cost $\mO(n^2)$ FLOPS
%         \end{itemize}
%         \bigskip
%     \begin{block}{Dense Matrix}
%         \begin{itemize}
%             \item A matrix is considered \textbf{dense} if the majority of its elements are non-zero. 
%             \item $A$ is dense $\implies$ $Ax=b$ takes more time to solve 
%         \end{itemize}        
%     \end{block}
%     \begin{alertblock}{Question}
%         How much does it cost to solve a dense matrix?
%     \end{alertblock}
% \end{frame}

%%%%%%%%%%%%%%%%%%%%%%%%%%%%%%%%%%%%%%%%%%%%%%%%%%%%%%%%%%%%%%%%%%%%%%%%

% \begin{frame}[fragile]
% \frametitle{Solving a general system}
% \begin{itemize}
%     \item Diagonal systems cost $\mO(n)$ FLOPs
%     \item Triangular systems cost $\mO(n^2)$ FLOPS
% \begin{alertblock}{Question}
%     What is the algorithm to solve a general system $Ax=b$?
% \end{alertblock}
% \bigskip
%     \item Formally the solution to $Ax=b$ is 
%  \begin{equation*}
%          x = A^{-1}b
%  \end{equation*}
%  \item \textbf{In MATLAB:}
%  \begin{lstlisting}[language=matlab]
%  >> A = ...
%  >> b = ...
%  >> x = A\b
%  \end{lstlisting}
% \end{itemize}
% \end{frame}

%%%%%%%%%%%%%%%%%%%%%%%%%%%%%%%%%%%%%%%%%%%%%%%%%%%%%%%%%%%%%%%%%%%%%%%%
\begin{frame}{Simple $2\times2$ Example}

\begin{example}
Solve
\begin{align*}
      x_1 + 3x_2 &= 5\\
    2 x_1 + 4x_2 &= 6
\end{align*}
\end{example}
\begin{itemize}
    \item Subtract 2 times the first equation from the second equation
\begin{align*}
      x_1 + 3x_2 &= 5\\
          - 2x_2 &= -4
\end{align*}

\onslide<2->{
   \item  The system is now in upper triangular form 
    \begin{equation*}
    \begin{bmatrix}1 & 3 \\ 0 & -2 \end{bmatrix}
    \left[\negthickspace\begin{array}{r}  x_1 \\ x_2 \end{array}\negthickspace\right] = 
    \left[\negthickspace\begin{array}{r}  5 \\ -4 \end{array}\negthickspace\right]
    \end{equation*}
    We can now use \textbf{backward substitution}.
}
\end{itemize}
 \end{frame}

%%%%%%%%%%%%%%%%%%%%%%%%%%%%%%%%%%%%%%%%%%%%%%%%%%%%%%%%%%%%%%%%%%%%%%%%
\begin{frame}{Simple $2\times2$ Example}


\begin{itemize}
    \item Solve the second equation for $x_2$
    \begin{equation*}
        x_2 = \frac{-4}{-2} = 2
    \end{equation*}
    \item Substitute  $x_2$ into the first equation and solve for $x_1$.
    \begin{equation*}
        x_1 = 5 - (3)(2) = -1
    \end{equation*}
\end{itemize}

\onslide<2>{

\keyidea{

    Change a system to triangular form and then do backward substitution!
    \begin{equation*}
    \begin{bmatrix}1 & 3 \\ 2 & 4 \end{bmatrix}
    \left[\negthickspace\begin{array}{r}  x_1 \\ x_2 \end{array}\negthickspace\right] = 
    \left[\negthickspace\begin{array}{r}  5 \\ 6 \end{array}\negthickspace\right]
    \rightarrow
    \begin{bmatrix}1 & 3 \\ 0 & -2 \end{bmatrix}
    \left[\negthickspace\begin{array}{r}  x_1 \\ x_2 \end{array}\negthickspace\right] = 
    \left[\negthickspace\begin{array}{r}  5 \\ -4 \end{array}\negthickspace\right]
    \rightarrow
    \text{back sub}
    \end{equation*}
}
}

\end{frame}

%%%%%%%%%%%%%%%%%%%%%%%%%%%%%%%%%%%%%%%%%%%%%%%%%%%%%%%%%%%%%%%%%%%%%%%%
\begin{frame}{General Procedure for $n\times n$}{Part 1}

\begin{itemize}
    \item[1.] Form the \textbf{augmented system} - matrix formed by appending $[A|b]$
    \[
    \begin{bmatrix}
        x & x & x & x & \big| & x \\
        x & x & x & x & \big| & x \\
        x & x & x & x & \big| & x \\
        x & x & x & x & \big| & x
    \end{bmatrix}
    \]

(The $x$'s represent numbers, they are generally \emph{not} the same
values.)\\
\vspace{0.5cm}
\item[2.] Choose your \textbf{pivot element} - the element which controls the multiplier. Start with the very first element $A_{11}$. The row containing the pivot element is called the \textbf{pivot row}.
    \[
    \begin{bmatrix}
        \boxed{x} & x & x & x  & x \\
        x & x & x & x &  x \\
        x & x & x & x &  x \\
        x & x & x & x &  x
    \end{bmatrix}
    \]

\end{itemize}
\end{frame}

%%%%%%%%%%%%%%%%%%%%%%%%%%%%%%%%%%%%%%%%%%%%%%%%%%%%%%%%%%%%%%%%%%%%%%%%
\begin{frame}[shrink]{General Procedure for $n\times n$}{Part 2}
\begin{itemize}
    \item[3.] Eliminate elements under the pivot element in the first column by performing \textbf{row operations} - multiply the pivot row by a scalar and add to a row below
\begin{center}
	\renewcommand{\arraystretch}{1.3}
	\footnotesize
	\begin{tabular}{cccc}
		$\begin{bmatrix}\boxed{x} & x & x & x & x \\
						   x  & x  & x  & x  & x  \\
						   x  & x  & x  & x  & x \\
						   x  & x  & x  & x  & x \end{bmatrix}$
		& $\longrightarrow$ &
		$\begin{bmatrix}\boxed{x} & x & x & x & x \\
						   0  & x' & x' & x' & x' \\
						   x  & x  & x  & x  & x \\
						   x  & x  & x  & x  & x \end{bmatrix}$ \\[10pt]
		& $\longrightarrow$ &
		$\begin{bmatrix}\boxed{x} & x & x & x & x \\
						   0  & x' & x' & x' & x' \\
						   0  & x' & x' & x' & x' \\
						   x  & x  & x  & x  & x \end{bmatrix}$ \\[10pt]
		& $\longrightarrow$ &
		$\begin{bmatrix}\boxed{x} & x & x & x & x \\
						   0  & x' & x' & x' & x' \\
						   0  & x' & x' & x' & x' \\
						   0  & x' & x' & x' & x'\end{bmatrix}$ &
	\end{tabular}
\end{center}
    ($x'$ indicates a value that has been changed once.)
\end{itemize}

\end{frame}
%%%%%%%%%%%%%%%%%%%%%%%%%%%%%%%%%%%%%%%%%%%%%%%%%%%%%%%%%%%%%%%%%%%%%%%%

\begin{frame}{General Procedure for $n\times n$}{Part 3}
\begin{itemize}
    \item[4.] Choose next pivot element: diagonal element in the second row.
    \item[5.] Eliminate elements under the pivot element in the second column.
\begin{center}
	\renewcommand{\arraystretch}{1.3}
	\footnotesize
	\begin{tabular}{ccccc}
		$\begin{bmatrix} x &   x        & x & x & x    \\
						 0 & \boxed{x'} & x' & x' & x' \\
						 0 &   x'       & x' & x' & x' \\
						 0 &   x'       & x' & x' & x' \end{bmatrix}$
		& $\longrightarrow$ &
		$\begin{bmatrix} x &   x        & x   & x   & x   \\
						 0 & \boxed{x'} & x'  & x'  & x'  \\
						 0 &   0        & x'' & x'' & x'' \\
						 0 &   x'       & x'  & x'  & x'\end{bmatrix}$ \\[30pt]
		& $\longrightarrow$ &
		$\begin{bmatrix} x &   x        & x   & x   & x   \\
						 0 & \boxed{x'} & x'  & x'  & x'  \\
						 0 &   0        & x'' & x'' & x'' \\
						 0 &   0        & x'' & x'' & x'' \end{bmatrix}$
	\end{tabular}
\end{center}
    ($x''$ indicates a value that has been changed twice.)
\end{itemize}

\end{frame}
%%%%%%%%%%%%%%%%%%%%%%%%%%%%%%%%%%%%%%%%%%%%%%%%%%%%%%%%%%%%%%%%%%%%%%%%
\begin{frame}{General Procedure for $n\times n$}{Part 4}
\begin{itemize}
    \item[6.] Repeat until the system is upper triangular.
\begin{center}
	\renewcommand{\arraystretch}{1.3}
	\footnotesize
	\begin{tabular}{ccccc}
		$\begin{bmatrix} x & x  &     x       & x   & x   \\
						 0 & x' &     x'      & x'  & x'  \\
						 0 & 0  & \boxed{x''} & x'' & x'' \\
						 0 & 0  &     x''     & x'' & x'' \end{bmatrix}$ 
		& $\longrightarrow$ &
		$\begin{bmatrix} x & x  &     x       & x    & x   \\
						 0 & x' &     x'      & x'   & x'  \\
						 0 & 0  & \boxed{x''} & x''  & x'' \\
						 0 & 0  &     0       & x''' & x''' \end{bmatrix}$ 
	\end{tabular}
\end{center}
    ($x'''$ indicates a value that has been changed three times.)
\end{itemize}

\vspace{3mm}

\onslide<2>{
\define{Gaussian Elimination}{
    The procedure of changing the augmented system to triangular form is called \textbf{Gaussian elimination} (GE).
}
}
\end{frame}
%%%%%%%%%%%%%%%%%%%%%%%%%%%%%%%%%%%%%%%%%%%%%%%%%%%%%%%%%%%%%%%%%%%%%%%%
\begin{frame}{Gaussian Elimination}

{\method{Gaussian Elimination}
{ $n\times n$ matrix $A$, $n \times 1$ right hand side $b$ }
{Upper triangular augmented system}
{This algorithm does not assume anything about $A$. It is a general solver. }}


\begin{block}{Getting the solution to $Ax=b$}
\begin{enumerate}

\item Construct the augmented system 

\item Perform Gaussian elimination
\begin{center}
	\renewcommand{\arraystretch}{1.3}
	\small
	\begin{tabular}{ccccc}
        $\begin{bmatrix}   x  & x & x & x & x \\
                           x  & x  & x  & x  & x  \\
                           x  & x  & x  & x  & x \\
                           x  & x  & x  & x  & x \end{bmatrix}$
		& $\longrightarrow$ &
		$\begin{bmatrix} x & x  &  x   & x    & x   \\
						 0 & x' &  x'  & x'   & x'  \\
						 0 & 0  &  x'' & x''  & x'' \\
						 0 & 0  &  0   & x''' & x''' \end{bmatrix}$ 
	\end{tabular}
\end{center}
\item Solve system with back sub 
\end{enumerate}
\end{block}
\end{frame}

%%%%%%%%%%%%%%%%%%%%%%%%%%%%%%%%%%%%%%%%%%%%%%%%%%%%%%%%%%%%%%%%%%%%%%%%
\begin{frame}{GE Example}{Part 1}

\begin{example}
Perform Gaussian elimination on the system $Ax=b$ where 
\begin{equation*}
    A = \left[\negthickspace\begin{array}{rrr} -3 &  2 & -1  \\
                                                6 & -6 &  7 \\
                                                3 & -4 &  4 \end{array}\negthickspace\right]
    \ \ \ \ \ \ \ \ \
    b = \left[\negthickspace\begin{array}{r} -1 \\ -7 \\ -6 \end{array}\negthickspace\right]
\end{equation*}
\end{example}
\begin{enumerate}
    \item[1.] Form the\textbf{ augmented system}
\begin{equation*}
    \tilde{A} = \left[ A\ \ b \right] =
      \left[
           \left.\negthickspace\begin{array}{rrr}  -3 &  2 & -1  \\
                                                    6 & -6 &  7 \\
                                                    3 & -4 &  4 \end{array}
           \right|\begin{array}{r} -1 \\ -7 \\ -6 \end{array}\negthickspace\right]
\end{equation*}
\end{enumerate}
\end{frame}

%%%%%%%%%%%%%%%%%%%%%%%%%%%%%%%%%%%%%%%%%%%%%%%%%%%%%%%%%%%%%%%%%%%%%%%%

\begin{frame}{GE Example}{Part 2}
\begin{enumerate}
    \item[2.] Choose 3 as the first pivot element. 
\begin{equation*}
    \tilde{A}  =
      \left[
           \left.\negthickspace\begin{array}{rrr}  \boxed{-3} &  2 & -1  \\
                                                    6 & -6 &  7 \\
                                                    3 & -4 &  4 \end{array}
           \right|\begin{array}{r} -1 \\ -7 \\ -6 \end{array}\negthickspace\right]
\end{equation*}

    \item[3.] How many times does $\boxed{-3}$ divide into $6$? Perform row operation to zero out the element below.
    \begin{itemize}
        \item $xmult =6/-3=-2$
        \item $R_2\leftarrow R_2-(xmult)R_1$
    \end{itemize}
\begin{equation*}
    \tilde{A}_{(1)} = \left[
           \left.\negthickspace\begin{array}{rrr}  -3 &  2 & -1 \\
                                                    0 & -2 &  5 \\
                                                    3 & -4 &  4 \end{array}
           \right|\begin{array}{r} -1 \\ -9 \\ -6 \end{array}\negthickspace\right]
\end{equation*}
\end{enumerate}
\end{frame}
%%%%%%%%%%%%%%%%%%%%%%%%%%%%%%%%%%%%%%%%%%%%%%%%%%%%%%%%%%%%
\begin{frame}{GE Example}{Part 3}
\begin{enumerate}

    \item[4.] How many times does $\boxed{-3}$ divide into $3$? Perform row operation to zero out the element below.
    \begin{itemize}
        \item $xmult =3/-3=-1$
        \item $R_3\leftarrow R_3-(xmult)R_1$
    \end{itemize}
    \begin{equation*}
    \tilde{A}_{(1)} = \left[
           \left.\negthickspace\begin{array}{rrr}  -3 &  2 & -1 \\
                                                    0 & -2 &  5 \\
                                                    3 & -4 &  4 \end{array}
           \right|\begin{array}{r} -1 \\ -9 \\ -6 \end{array}\negthickspace\right]
\end{equation*}
\begin{equation*}
    \tilde{A}_{(1)} = \left[
           \left.\negthickspace\begin{array}{rrr}  -3 &  2 & -1 \\
                                                    0 & -2 &  5 \\
                                                    0 & -2 &  3 \end{array}
           \right|\begin{array}{r} -1 \\ -9 \\ -7 \end{array}\negthickspace\right]
\end{equation*}
\end{enumerate}
\end{frame}
%%%%%%%%%%%%%%%%%%%%%%%%%%%%%%%%%%%%%%%%%%%%%%%%%%%%%%%%%%%%

\begin{frame}{GE Example}{Part 4}
\begin{enumerate}
    \item[5.] Choose element -2 as the next pivot 
    \begin{equation*}
        \tilde{A}_{(1)} = \left[
               \left.\negthickspace\begin{array}{rrr}  -3 &  2 & -1 \\
                                                        0 & \boxed{-2} &  5 \\
                                                        0 & -2 &  3 \end{array}
               \right|\begin{array}{r} -1 \\ -9 \\ -7 \end{array}\negthickspace\right]
    \end{equation*}

    \item[6.] How many times does $\boxed{-2}$ divide into $-2$? Perform row operation to zero out the element below.
    \begin{itemize}
        \item $xmult = -2/-2=1$
        \item $R_3\leftarrow R_3 -(xmult)R_2$
    \end{itemize}

\begin{equation*}
    \tilde{A}_{(2)} = \left[
           \left.\negthickspace\begin{array}{rrr}  -3 &  2 & -1 \\
                                                    0 & -2 &  5 \\
                                                    0 &  0 & -2 \end{array}
           \right|\begin{array}{r} -1 \\ -9 \\ 2 \end{array}\negthickspace\right]
\end{equation*}
\end{enumerate}
\vfill
\begin{center}
    Done.
\end{center}
\end{frame}
%%%%%%%%%%%%%%%%%%%%%%%%%%%%%%%%%%%%%%%%%%%%%%%%%%%%%%%%%%%%%%%%%%%%%%%%

\begin{frame}{Solving the System}

\cols{ 
\begin{block}{Option 1: Back Sub}   
\begin{equation*}
    \tilde{A}_{(2)} = \left[
           \left.\negthickspace\begin{array}{rrr}  -3 &  2 & -1 \\
                                                    0 & -2 &  5 \\
                                                    0 &  0 & -2 \end{array}
           \right|\begin{array}{r} -1 \\ -9 \\ 2 \end{array}\negthickspace\right]
\end{equation*}
Solve by back substitution to get
\begin{align*}
    x_3 &= \frac{\ \,2}{-2} = -1                             \\
    x_2 &= \frac{\ \,1}{-2}\left(-9 -5x_3 \right) = 2        \\
    x_1 &= \frac{\ \,1}{-3}\left(-1 -2x_2 + x_3 \right) = 2  \\
\end{align*}
\end{block}
}{

\begin{block}{Option 2: Gauss-Jordan}   
    Continue doing operations to transform matrix to an identity
\begin{equation*}
        \left[
           \left.\negthickspace\begin{array}{rrr}  -3 &  2 & -1 \\
                                                    0 & -2 &  5 \\
                                                    0 &  0 & \boxed{1} \end{array}
           \right|\begin{array}{r} -1 \\ -9 \\ -1 \end{array}\negthickspace\right]
\end{equation*}
\begin{equation*}
           \left[
           \left.\negthickspace\begin{array}{rrr}  -3 &  2 &  0 \\
                                                    0 & -2 &  0 \\
                                                    0 &  0 & 1 \end{array}
           \right|\begin{array}{r} -2 \\ -4 \\ -1 \end{array}\negthickspace\right]
\end{equation*}
           $$\vdots$$
\begin{equation*}
           \left[
           \left.\negthickspace\begin{array}{rrr}  1 &  0 &  0 \\
                                                    0 & 1 &  0 \\
                                                    0 &  0 & 1 \end{array}
           \right|\begin{array}{r} 2 \\ 2 \\ -1 \end{array}\negthickspace\right]
\end{equation*}
\end{block}
}
\end{frame}

%%%%%%%%%%%%%%%%%%%%%%%%%%%%%%%%%%%%%%%%%%%%%%%%%%%%%%%%%%%%%%%%%%%%%%%%

\begin{frame}[fragile]{Gaussian Elimination Algorithm}

\begin{lstlisting}[mathescape,label=algo:fwdelimb]
given $A$, $b$
                           
for $k=1\dots n-1$          
    for $i = k+1\ldots n$      
        for $j = k\ldots n$      
            $a_{ij} = a_{ij} - (a_{ik}/a_{kk}) a_{kj}$
        end                      
        $b_{i} = b_{i} - (a_{ik}/a_{kk}) b_{k}$
    end                        
end                          
\end{lstlisting}
\vspace{0.3cm}

\begin{block}{Two ways to Optimize}
    
\begin{itemize}
    \item[1.] Precompute the multiplier $a_{ik}/a_{kk}$ before line 5
    \item[2.] Line 6, if $j=k$ the value $a_{ij}$ is always zero. Just skip it.
\end{itemize}
\end{block}



\end{frame}

%%%%%%%%%%%%%%%%%%%%%%%%%%%%%%%%%%%%%%%%%%%%%%%%%%%%%%%%%%%%%%%%%%%%%%%%
\begin{frame}[fragile]{Gaussian Elimination Algorithm}
%      $a_{ik} = xmult$
\begin{lstlisting}[mathescape,label=algo:fwdelim,caption=Optimized Version]
given $A$, $b$
                               
for $k=1\dots n-1$          
    for $i = k+1\ldots n$      
        $xmult = a_{ik}/a_{kk}$     % precompute multiplier
        $a_{ik} = 0$                % zero out elements below pivot
        for $j = k+1\ldots n$        % start loop at j=k+1
            $a_{ij} = a_{ij} - (xmult) a_{kj}$
        end                      
        $b_{i} = b_{i} - (xmult) b_{k}$
    end                        
end                          
\end{lstlisting}
\vspace{3mm}

\end{frame}

%%%%%%%%%%%%%%%%%%%%%%%%%%%%%%%%%%%%%%%%%%%%%%%%%%%%%%%%%%%%%%%%%%%%%%
\begin{frame}[fragile]{Gaussian Elimination Algorithm: Storing Multipliers}
\begin{lstlisting}[mathescape,label=algo:fwdelim]
given $A$, $b$
                           
for $k=1\dots n-1$          
    for $i = k+1\ldots n$      
        $xmult = a_{ik}/a_{kk}$
        $a_{ik} = xmult$      % save multiplier
        for $j = k+1\ldots n$      
            $a_{ij} = a_{ij} - (xmult) a_{kj}$
        end                      
        $b_{i} = b_{i} - (xmult) b_{k}$
    end                        
end                          
\end{lstlisting}
\vspace{3mm}

\begin{itemize}
    \item This variant \textit{saves} the multipliers by writing them into the matrix 
    \item We will discuss why this might be desirable 
\end{itemize}

\end{frame}
%%%%%%%%%%%%%%%%%%%%%%%%%%%%%%%%%%%%%%%%%%%%%%%%%%%%%%%%%%%%%%%%%%%%%%%%
\begin{frame}[fragile]{Zeroing vs. Storing Multipliers}

\begin{itemize}
    \item Zeroing out version 
    $$
      \left[
           \left.\negthickspace\begin{array}{rrr}  -3 &  2 & -1  \\
                                                    6 & -6 &  7 \\
                                                    3 & -4 &  4 \end{array}
           \right|\begin{array}{r} -1 \\ -7 \\ -6 \end{array}\negthickspace\right]\rightarrow \left[
           \left.\negthickspace\begin{array}{rrr}  -3 &  2 & -1 \\
                                                    0 & -2 &  5 \\
                                                    0 &  0 & -2 \end{array}
           \right|\begin{array}{r} -1 \\ -9 \\ 2 \end{array}\negthickspace\right]
    $$
    \item Storing the multipliers $xmult$ in the lower triangle
    $$
      \left[
           \left.\negthickspace\begin{array}{rrr}  -3 &  2 & -1  \\
                                                    6 & -6 &  7 \\
                                                    3 & -4 &  4 \end{array}
           \right|\begin{array}{r} -1 \\ -7 \\ -6 \end{array}\negthickspace\right]\rightarrow \left[
           \left.\negthickspace\begin{array}{rrr}  -3 &  2 & -1 \\
                                                    \textcolor{red}{-2} & -2 &  5 \\
                                                    \textcolor{red}{-1} &  \textcolor{red}{1} & -2 \end{array}
           \right|\begin{array}{r} -1 \\ -9 \\ 2 \end{array}\negthickspace\right]
    $$
\end{itemize}

\question{Why would you store multipliers?}

\end{frame}

%%%%%%%%%%%%%%%%%%%%%%%%%%%%%%%%%%%%%%%%%%%%%%%%%%%%%%%%%%%%%%%%%%%%%%%%
\sectionframe{Cost of GE}
%%%%%%%%%%%%%%%%%%%%%%%%%%%%%%%%%%%%%%%%%%%%%%%%%%%%%%%%%%%%%%%%%%%%%%%%
\begin{frame}[fragile]{Naive Gaussian Elimination Algorithm}

\question{
    What is the cost of Gaussian elimination? \\[5pt]
    \quad or rather\\[5pt]
    What is the cost to solve a \underline{general system}?
}


\vspace{5mm}
\begin{itemize}
    \item Solve = Gaussian Elimination + Back Substitution
    \item We already know Backward Substitution is $\mathcal{O}(n^2)$
    \item What about Gaussian elimination?
\end{itemize}
\vfill

% \begin{block}{GE\_naive.m $\quad$
% GE\_naive\_test.m}
% Since we have the algorithm, we can experimentally benchmark GE. 
% \end{block}
% ----------------------
\end{frame}
%%%%%%%%%%%%%%%%%%%%%%%%%%%%%%%%%%%%%%%%%%%%%%%%%%%%%%%%%%%%%%%%%%%%%%%%
\begin{frame}{Cost for a $3\times 3$}{Part 1}


\question{
What is the exact number of FLOPs for the $3\times 3$ example?
\begin{equation*}
\small
      \left[\left.\negthickspace\begin{array}{rrr}  -3 &  2 & -1  \\
                                                    6 & -6 &  7 \\
                                                    3 & -4 &  4 \end{array}
       \right|\begin{array}{r} -1 \\ -7 \\ -6 \end{array}\negthickspace\right]
            \longrightarrow
            \left[\left.\negthickspace\begin{array}{rrr}  -3 &  2 & -1 \\
                                                    0 & -2 &  5 \\
                                                    0 &  0 & -2 \end{array}
           \right|\begin{array}{r} -1 \\ -9 \\ 2 \end{array}\negthickspace\right]
\end{equation*}
}

\begin{itemize}
\item Step 1: Eliminate $a_{21}$
$$
      \left[\left.\negthickspace\begin{array}{rrr}  \boxed{-3} &  2 & -1  \\
                                                    \red{6} & -6 &  7 \\
                                                    3 & -4 &  4 \end{array}
       \right|\begin{array}{r} -1 \\ -7 \\ -6 \end{array}\negthickspace\right]
            \longrightarrow
        \left[\left.\negthickspace\begin{array}{rrr}  -3 &  2 & -1 \\
                                                    0 & -2 &  5 \\
                                                    3 & -4 &  4 \end{array}
           \right|\begin{array}{r} -1 \\ -9 \\ -6 \end{array}\negthickspace\right]
$$
\begin{itemize}
    \item 1 division: $6/-3 = -2$
    \item 3 multiplications: $-2R_1 = [\phantom{-}6 \phantom{-}-4\; \phantom{-}2\; | \phantom{-}2]$  
    \item 3 subtractions: $R_2 - 2 R_1 = [\phantom{-}0\; -2\; \phantom{-}5\; | -9]$ 
\end{itemize}
\end{itemize}
\begin{alertblock}{}
Note: We don't count $6-(-3)(-2)=0$ at $A_{21}$ because the result it always $0$
\end{alertblock}
\end{frame}
%%%%%%%%%%%%%%%%%%%%%%%%%%%%%%%%%%%%%%%%%%%%%%%%%%%%%%%%%%%%%%%%%%%%%%%%
\begin{frame}{Cost for a $3\times 3$}{Part 2}
\begin{itemize}

\item Step 2: Eliminate $a_{31}$
$$
        \left[\left.\negthickspace\begin{array}{rrr} \boxed{-3} &  2 & -1 \\
                                                    0 & -2 &  5 \\
                                                    \red{3} & -4 &  4 \end{array}
           \right|\begin{array}{r} -1 \\ -9 \\ -6 \end{array}\negthickspace\right]
            \longrightarrow
        \left[\left.\negthickspace\begin{array}{rrr}  -3 &  2 & -1 \\
                                                    0 & -2 &  5 \\
                                                    0 & -2 &  3 \end{array}
           \right|\begin{array}{r} -1 \\ -9 \\ -7 \end{array}\negthickspace\right]
$$
\begin{itemize}
    \item 1 division: $3/-3 = -1$
    \item 3 multiplications: $-R_1 = [\phantom{-}3\; -2\; \phantom{-}1\; | \phantom{-}1]$
    \item 3 subtractions: $R_3 - (-R_1) = [\phantom{-}0\; -2\; \phantom{-}3\; | -7]$ 
\end{itemize}

\item Step 3: Eliminate $a_{23}$
$$
    \left[\left.\negthickspace\begin{array}{rrr}  -3 &  2 & -1 \\
                                                0 & \boxed{-2} &  5 \\
                                                0 & \red{-2} &  3 \end{array}
       \right|\begin{array}{r} -1 \\ -9 \\ -7 \end{array}\negthickspace\right]
        \longrightarrow
        \left[\left.\negthickspace\begin{array}{rrr}  -3 &  2 & -1 \\
                                                0 & -2 &  5 \\
                                                0 &  0 & -2 \end{array}
       \right|\begin{array}{r} -1 \\ -9 \\ 2 \end{array}\negthickspace\right]
$$
\begin{itemize}
    \item 1 division: $-2/-2 = 1$
    \item 2 multiplications: $ (1)R_1 = [\phantom{-}0\; -2\; \phantom{-}5\; | -9]$
    \item 2 subtractions: $R_3 - R_1 = [\phantom{-}0\; \phantom{-}0\; -2\; | \phantom{-}2]$ 
\end{itemize}
\end{itemize}
\end{frame}
%%%%%%%%%%%%%%%%%%%%%%%%%%%%%%%%%%%%%%%%%%%%%%%%%%%%%%%%%%%%%%%%%%%%%%%%
\begin{frame}{Cost for a $3\times 3$}{Part 3}
\begin{itemize}

\vspace{5mm}
\item FLOP summary
\begin{center}
\begin{tabular}{c c c c c c}\hline
 Row & Step & Add & Multiply & Divide & Total \\\hline
 1 &  Eliminate element $a_{21}$ & 3 & 3& 1 & 7\\
 1 &  Eliminate element $a_{31}$ & 3 & 3& 1 & 7\\
 2 &  Eliminate element $a_{32}$ & 2 & 2 & 1 & 5\\\hline
\end{tabular}
\end{center}
\vspace{5mm}
\item Total cost: 19 FLOPs

\begin{block}{What's the pattern?}
Let $n=3$. 
\begin{itemize}
    \item In row 1, there are $n-1=2$ elements that need to be eliminated. Each elimination requires $n$ additions, $n$ multiplications, and 1 division. 
    \item In row 2, there are $n-2=1$ element(s) that need to be eliminated. Each elimination requires $n-1$ additions, $n-1$ multiplications, and 1 division.
    
\end{itemize}
\end{block}
\end{itemize}
\end{frame}
%%%%%%%%%%%%%%%%%%%%%%%%%%%%%%%%%%%%%%%%%%%%%%%%%%%%%%%%%%%%%%%%%%%%%%%%
\begin{frame}[shrink]{What about a $4\times4$?}

\begin{itemize}
    \item We have 6 elimination targets: 
    $
    \begin{bmatrix}
        x & x & x & x  \\
        \red{x} & x & x & x  \\
        \red{x} & \red{x}  & x & x  \\
        \red{x} & \red{x}  & \red{x}  & x 
    \end{bmatrix}
    $
    
    \item This is the pattern for a $3\times 3$
\begin{center}
\begin{tabular}{c c c c c c}\hline
 Row & Step & Add & Multiply & Divide & Total \\\hline
 1 &  Eliminate element $a_{21}$ & 3 & 3& 1 & 7\\
 1 &  Eliminate element $a_{31}$ & 3 & 3& 1 & 7\\
 2 &  Eliminate element $a_{32}$ & 2 & 2 & 1 & 5\\\hline
\end{tabular}
\vspace{1em}
\end{center}
\item<2>  This is the pattern for a $4\times 4$
\begin{center}
\begin{tabular}{c c c c c c}\hline
 Row & Step & Add & Multiply & Divide & Total \\\hline
 1 &  Eliminate element $a_{21}$ & 4 & 4& 1 & 9\\
 1 &  Eliminate element $a_{31}$ & 4 & 4& 1 & 9\\
 1 &  Eliminate element $a_{41}$ & 4 & 4& 1 & 9\\
 2 &  Eliminate element $a_{32}$ & 3 & 3 & 1 & 7\\
 2 &  Eliminate element $a_{42}$ & 3 & 3 & 1 & 7\\
 3 &  Eliminate element $a_{34}$ & 2 & 2 & 1 & 5\\\hline
\end{tabular}
\end{center}
\end{itemize}
\end{frame}
%%%%%%%%%%%%%%%%%%%%%%%%%%%%%%%%%%%%%%%%%%%%%%%%%%%%%%%%%%%%%%%
\begin{frame}{The Pattern for $4\times4$}
    \begin{center}
\begin{tabular}{c c c c c c}\hline
 Row & Step & Add & Multiply  \\\hline
 1 &  Eliminate element $a_{21}$ & 4 & 4\\
 1 &  Eliminate element $a_{31}$ & 4 & 4\\
 1 &  Eliminate element $a_{41}$ & 4 & 4\\
 2 &  Eliminate element $a_{32}$ & 3 & 3 \\
 2 &  Eliminate element $a_{42}$ & 3 & 3\\
 3 &  Eliminate element $a_{34}$ & 2 & 2 \\\hline
\end{tabular}
\end{center}
Let $n=4$
\begin{itemize}
    \item \textbf{Row 1:} $4+4+4=n+n+n=(n-1)n$ Adds 
    \item \textbf{Row 2:} $3+3=(n-1)+(n-1)=(n-2)(n-1)$ Adds
    \item \textbf{Row 3:} $2=(n-2)=(n-3)(n-2)$ Adds
\end{itemize}
\begin{alertblock}{}
    There is a consistent pattern in the way the number of operations decrease.
\end{alertblock}
\end{frame}
%%%%%%%%%%%%%%%%%%%%%%%%%%%%%%%%%%%%%%%%%%%%%%%%%%%%%%%%%%%%%%%
\begin{frame}{How many additions?}

\question{How many total additions are required for $4\times 4$ system?}

$$
\underbrace{(n-1)n}_{row\,1} + \underbrace{(n-2)(n-1)}_{row\,2}+\underbrace{(n-3)(n-2)}_{row \,3}=\sum_{j=2}^4(j-1)j
$$
\onslide<2>{
\question{How many total additions are required for $n\times n$ system?}
$$
\sum_{j=2}^n j(j-1) \text{ additions}
$$
}


\end{frame}
%%%%%%%%%%%%%%%%%%%%%%%%%%%%%%%%%%%%%%%%%%%%%%%%%%%%%%%%%%%%%%%
\begin{frame}{Cost for a $n\times n$}{Part 1}
\begin{itemize}
    \item Next, let's extrapolate the results for the $3 \times 3$ to determine cost for a general $n \times n$.\\[3mm]

    \item Key is to count operations by rows\\[3mm]


    \item Operation count:
    
\begin{center}
\begin{tabular}{c c c c}\hline
Row $j$ & Add & Multiply & Divide\\\hline
1 & $n(n-1)$ & $n(n-1)$ & $n-1$\\
2 & $(n-1)(n-2)$ & $(n-1)(n-2)$ & $n-2$\\
\vdots & \vdots &\vdots &\vdots \\
$n-1$ & $2$ & $2$ & $1$\\\hline
Total: &  $\sum_{j=2}^{n} j(j-1)$ & $\sum_{j=2}^{n} j(j-1)$ &  $\sum_{j=1}^{n-1} j$
\end{tabular}
\end{center}

\end{itemize}
\end{frame}
%%%%%%%%%%%%%%%%%%%%%%%%%%%%%%%%%%%%%%%%%%%%%%%%%%%%%%%%%%%%%%%%%%%%%%%%
\begin{frame}{Cost for a $n\times n$}{Part 2}
\begin{itemize}
    \item Total operations as summations
\begin{center}
\begin{tabular}{c c}\hline\\
add & $\sum_{j=2}^{n} j(j-1) \approx \sum_{j=1}^nj^2$ \\[5pt]
multiply & $\sum_{j=2}^{n} j(j-1) \approx \sum_{j=1}^nj^2 $\\[5pt]
divide & $\sum_{j=1}^{n-1} j$ \\[7pt]
\hline
\end{tabular}
\end{center}

\vspace{5mm}

\item Let's use the formulas $\sum_{j=1}^{p} j = \frac{p(p+1)}{2}$ and
$\sum_{j=1}^{p} j^2= \frac{p(p+1)(2p+1)}{6}$
\begin{center}
\begin{tabular}{c c}\hline\\
add & $\frac{n(n-1)(2n-1)}{6}$ \\[8pt]
multiply & $\frac{n(n-1)(2n-1)}{6}$\\[8pt]
divide & $\frac{n(n-1)}{2}$ \\[7pt] 
\hline
\end{tabular}
\end{center}
\end{itemize}
\end{frame}

%%%%%%%%%%%%%%%%%%%%%%%%%%%%%%%%%%%%%%%%%%%%%%%%%%%%%%%%%%%%%%%%%%%%%%%%

\begin{frame}{Cost for a $n\times n$}{Part 3}
\begin{itemize}
    \item Solve = Gaussian Elimination + Back Substitution  
    \item \textbf{Gaussian elimination}
    \begin{center}
\begin{tabular}{c c}\hline\\
add & $\frac{n(n-1)(2n-1)}{6}$ \\[4pt]
multiply & $\frac{n(n-1)(2n-1)}{6}$\\[4pt]
divide & $\frac{n(n-1)}{2}$ \\[4pt] 
\hline
\end{tabular}
\end{center}
    \item \textbf{Back substitution}
\begin{center}
\begin{tabular}{c c}\hline \\
add-subtracts & $\frac{n(n - 1)}{2}$\\[4pt]
multiply-divides & $\frac{n(n + 1)}{2}$\\[4pt]
\hline
\end{tabular}
\end{center}
\end{itemize}
\end{frame}
%%%%%%%%%%%%%%%%%%%%%%%%%%%%%%%%%%%%%%%%%%%%%%%%%%%%%%%%%%%%%%%%%%%%%%%%
\begin{frame}{Total Cost for Naive Gaussian Elimination}
Combining the cost of forward elimination, updating $b$, and backward substitution
gives
\begin{center}
\begin{tabular}{c c}\hline\\
add-subtracts & $\frac{n(n-1)(2n-1)}{6} + \frac{n(n - 1)}{2} + \frac{n(n - 1)}{2}$\\[5pt]
&  $=\frac{n(n-1)(2n+5)}{6}= \frac{1}{3}n^3 + \mO(n^2)$\\[6pt]
multiply-divides & $\frac{n(n^2 - 1)}{3}+\frac{n(n - 1)}{2}+ \frac{n(n + 1)}{2}$\\[5pt]
& $=\frac{n(n^2 + 3n-1)}{3}=\frac{1}{3}n^3 + \mO(n^2)$\\[6pt]
\hline
\end{tabular}
\end{center}

\keyidea{
\begin{itemize}
    \item The cost to solve a general system $Ax=b$ is
\begin{equation*}
  \frac{2}{3} n^3  + \mO(n^2) \quad \textrm{FLOPs}
\end{equation*}
    \item $\mO(n^3)$ is expensive but it is necessary 
\end{itemize}
}
\end{frame}

%%%%%%%%%%%%%%%%%%%%%%%%%%%%%%%%%%%%%%%%%%%%%%%%%%%%%%%%%%%%%%%%%%%%%%%%
\sectionframe{Why ``Naive''?}
%%%%%%%%%%%%%%%%%%%%%%%%%%%%%%%%%%%%%%%%%%%%%%%%%%%%%%%%%%%%%%%%%%%%%%%%
\begin{frame}{}

\question{
Why is the algorithm called the ``Naive'' Gaussian elimination?
}
\vfill

\begin{itemize}
    \item  It is called naive because of the simplistic \textbf{choice of pivot elements }
    \[
    \begin{bmatrix}
        \boxed{x} & x & x & x  & x \\
        x & \boxed{x} & x & x &  x \\
        x & x & \boxed{x} & x &  x \\
        x & x & x & x &  x
    \end{bmatrix}
    \]
    \item Naive GE choose the pivot elements to be $a_{ii}$ for $i = 1$ to $n-1$
    \vfill 
    \item There are \textit{drawbacks} to this method
    \vfill 
\end{itemize}

    
% \begin{equation*}
% A=
% \begin{bmatrix}
%   0 & 2 & 3\\
%   4 & 5 & 6\\
%   7 & 8 & 9\\
% \end{bmatrix}
% \end{equation*}

% \bigskip
% \begin{equation*}
% B=
% \begin{bmatrix}
%   10^{-10} & 2 & 3\\
%   4 & 5 & 6\\
%   7 & 8 & 9\\
% \end{bmatrix}
% \end{equation*}
\end{frame}

%%%%%%%%%%%%%%%%%%%%%%%%%%%%%%%%%%%%%%%%%%%%%%%%%%%%%%%%%%%%%%%%%%%%%%%%

\begin{frame}{Zero on $a_{11}$}

\begin{itemize}
    \item Consider performing GE on this matrix
\begin{equation*}
A=
\begin{bmatrix}
  0 & 2 & 3\\
  4 & 5 & 6\\
  7 & 8 & 9\\
\end{bmatrix}
\end{equation*}
    \item Naively choosing the pivot $a_{11} = 0$ results in division by zero 
\begin{equation*}
A=
\begin{bmatrix}
  \boxed{0} & 2 & 3\\
  4 & 5 & 6\\
  7 & 8 & 9\\
\end{bmatrix}
\end{equation*}
$$
R_2 \leftarrow R_2 - \frac{4}{\red{0}} R_1
$$
\item Naive GE will not work whenever $a_{11} = 0$
\end{itemize}
\end{frame}

%%%%%%%%%%%%%%%%%%%%%%%%%%%%%%%%%%%%%%%%%%%%%%%%%%%%%%%%%%%%%%%%%%%%%%%%

\begin{frame}{Another Example}
\begin{itemize}
    \item Try Naive GE on
        \begin{equation*}
            A = \left[\negthickspace\begin{array}{rrrr}  2 &  4 & -2  \\
                                                         1 &  2 &  4   \\
                                                        -2 & -3 &  8   \end{array}\negthickspace\right]
            \ \ \ \ \ \ \ \ \
            b = \left[\negthickspace\begin{array}{r} -4 \\  5 \\ 7 \end{array}\negthickspace\right]
        \end{equation*}
        \item Choose $a_{11}=2$ as the pivot element
\begin{equation*}
%   \tilde{A} =
    \left[
           \left.\negthickspace\begin{array}{rrrr}  \boxed{2} &  4 & -2  \\
                                                    1 &  2 &  4  \\
                                                   -2 & -3 &  8  \end{array}
           \right|\begin{array}{r} -4 \\  5 \\ 7  \end{array}\negthickspace\right]
           \quad \begin{aligned}
                       &R_2 \leftarrow-\frac{1}{2}R_1 + R_2 \\
         &R_3 \leftarrow R_1 + R_3
           \end{aligned}
\end{equation*}
    \item Result: 
        \begin{equation*}
        \left[\left.\negthickspace\begin{array}{rrrr}  2 &  4 & -2 \\
                                                        0 &  \boxed{0} &  5 \\
                                                        0 &  1 &  6  \end{array}
               \right|\begin{array}{r} -4 \\ 7 \\ -3 \end{array}\negthickspace\right]
    \end{equation*}
    \item \textbf{Problem!} The next pivot element $a_{22}$  is a zero!
    \end{itemize}
\end{frame}

%%%%%%%%%%%%%%%%%%%%%%%%%%%%%%%%%%%%%%%%%%%%%%%%%%%%%%%%%%%%%%%%%

\begin{frame}{Small Pivots}
\begin{itemize}
    \item But this is \textbf{even worse}
    \begin{equation*}
A=
\begin{bmatrix}
  \boxed{10^{-16}} & 2 & 3\\
  4 & 5 & 6\\
  7 & 8 & 9\\
\end{bmatrix}
\end{equation*}
    \item<2> After one row operation... 
        \begin{equation*}
    \begin{bmatrix}
  10^{-16} & 2 & 3\\
  0 & ??? & ???\\
  7 & 8 & 9\\
\end{bmatrix}
\end{equation*}
    \question{What happens to the $5$ and $6$ in the second row?}
\end{itemize}
\end{frame}
%%%%%%%%%%%%%%%%%%%%%%%%%%%%%%%%%%%%%%%%%%%%%%%%%%%%%%%%%%%%%%%%%
\begin{frame}{What's Next?}
\begin{itemize}
\item How to deal with a zero pivot  
\vfill

\item Three types of pivoting algorithms: 
\begin{enumerate}
    \item Partial Pivoting
    \item Full Pivoting 
    \item Scaled Partial Pivoting
\end{enumerate}
\vfill

\item Cost of each pivoting algorithm
\vfill

\item Effect of small pivot elements in rounding
\vfill
\end{itemize}
\end{frame}

\end{document}
